\begin{lemma}
\label{lemma-crystals-on-affine}
In the situation above there is a functor
$$
\begin{matrix}
\text{crystals in quasi-coherent} \\
\mathcal{O}_{X/S}\text{-modules on }\text{Cris}(X/S)
\end{matrix}
\longrightarrow
\begin{matrix}
\text{pairs }(M, \nabla)\text{ satisfying} \\
\text{(\ref{item-complete}), (\ref{item-connection}),
(\ref{item-integrable}), and (\ref{item-topologically-quasi-nilpotent})}
\end{matrix}
$$
\end{lemma}

\begin{proof}
Let $\mathcal{F}$ be a crystal in quasi-coherent modules on $X/S$.
Set $T_e = \Spec(D_e)$ so that $(X, T_e, \bar\gamma)$ is an object
of $\text{Cris}(X/S)$ for $e \gg 0$. We have morphisms
$$
(X, T_e, \bar\gamma) \to (X, T_{e + 1}, \bar\gamma) \to \ldots
$$
which are closed immersions. We set
$$
M =
\lim_e \Gamma((X, T_e, \bar\gamma), \mathcal{F}) =
\lim_e \Gamma(T_e, \mathcal{F}_{T_e}) = \lim_e M_e
$$
Note that since $\mathcal{F}$ is locally quasi-coherent we have
$\mathcal{F}_{T_e} = \widetilde{M_e}$. Since $\mathcal{F}$ is a
crystal we have $M_e = M_{e + 1}/p^eM_{e + 1}$. Hence we see that
$M_e = M/p^eM$ and that $M$ is $p$-adically complete, see
Algebra, Lemma \ref{algebra-lemma-limit-complete}.

\medskip\noindent
By Lemma \ref{lemma-automatic-connection} we know that $\mathcal{F}$
comes endowed with a canonical integrable connection
$\nabla : \mathcal{F} \to \mathcal{F} \otimes \Omega_{X/S}$.
If we evaluate this connection on the objects $T_e$ constructed above
we obtain a canonical integrable connection
$$
\nabla : M \longrightarrow M \otimes^\wedge_D \Omega_D
$$
To see that this is topologically nilpotent we work out what this means.

\medskip\noindent
Now we can do the same procedure for the rings $D(n)$.
This produces a $p$-adically complete $D(n)$-module $M(n)$. Again using
the crystal property of $\mathcal{F}$ we obtain isomorphisms
$$
M \otimes^\wedge_{D, p_0} D(1) \rightarrow M(1)
\leftarrow M \otimes^\wedge_{D, p_1} D(1)
$$
compare with the proof of Lemma \ref{lemma-automatic-connection}.
Denote $c$ the composition from left to right. Pick $m \in M$.
Write $\xi_i = x_i \otimes 1 - 1 \otimes x_i$.
Using (\ref{equation-structure-Dn}) we can write uniquely
$$
c(m \otimes 1) = \sum\nolimits_K \theta_K(m) \otimes \prod \xi_i^{[k_i]}
$$
for some $\theta_K(m) \in M$ where the sum is over multi-indices
$K = (k_i)$ with $k_i \geq 0$ and $\sum k_i < \infty$. Set
$\theta_i = \theta_K$ where $K$ has a $1$ in the $i$th spot and
zeros elsewhere. We have
$$
\nabla(m) = \sum \theta_i(m) \text{d}x_i.
$$
as can be seen by comparing with the definition of
$\nabla$. Namely, the defining equation is
$p_1^*m = \nabla(m) - c(p_0^*m)$ in Lemma \ref{lemma-automatic-connection}
but the sign works out because in the Stacks project we consistently use
$\text{d}f = p_1(f) - p_0(f)$ modulo the ideal of the diagonal squared,
and hence $\xi_i = x_i \otimes 1 - 1 \otimes x_i$ maps to $-\text{d}x_i$
modulo the ideal of the diagonal squared.

\medskip\noindent
Denote $q_i : D \to D(2)$ and $q_{ij} : D(1) \to D(2)$ the coprojections
corresponding to the indices $i, j$. As in the last paragraph of the proof of
Lemma \ref{lemma-automatic-connection}
we see that
$$
q_{02}^*c = q_{12}^*c \circ q_{01}^*c.
$$
This means that
$$
\sum\nolimits_{K''} \theta_{K''}(m) \otimes \prod {\zeta''_i}^{[k''_i]}
=
\sum\nolimits_{K', K} \theta_{K'}(\theta_K(m))
\otimes \prod {\zeta'_i}^{[k'_i]} \prod \zeta_i^{[k_i]}
$$
in $M \otimes^\wedge_{D, q_2} D(2)$ where
\begin{align*}
\zeta_i & = x_i \otimes 1 \otimes 1 - 1 \otimes x_i \otimes 1,\\
\zeta'_i & = 1 \otimes x_i \otimes 1 - 1 \otimes 1 \otimes x_i,\\
\zeta''_i & = x_i \otimes 1 \otimes 1 - 1 \otimes 1 \otimes x_i.
\end{align*}
In particular $\zeta''_i = \zeta_i + \zeta'_i$ and we have that
$D(2)$ is the $p$-adic completion of the divided power polynomial
ring in $\zeta_i, \zeta'_i$ over $q_2(D)$, see Lemma \ref{lemma-structure-Dn}.
Comparing coefficients in the expression above it follows immediately that
$\theta_i \circ \theta_j = \theta_j \circ \theta_i$
(this provides an alternative proof of the integrability of $\nabla$) and that
$$
\theta_K(m) = (\prod \theta_i^{k_i})(m).
$$
In particular, as the sum expressing $c(m \otimes 1)$ above has to converge
$p$-adically we conclude that for each $i$ and each $m \in M$ only a finite
number of $\theta_i^k(m)$ are allowed to be nonzero modulo $p$.
\end{proof}